% Part: applied-modal-logics
% Chapter: temporal-logic
% Section: language-epistemic-logic

\documentclass[../../../include/open-logic-section]{subfiles}

\begin{document}

\olfileid{aml}{tl}{sem}

\olsection{د زماني منطق معنٰی پوهنه}

\begin{defn}
د زماني منطق بنسټيزه ژبه دا توکي لري:
\begin{enumerate}
  \tagitem{prvFalse}{د !!{falsity} قضيوي ثابت~$\lfalse$۔}{}
  \tagitem{prvTrue}{د !!{truth} قضيوي ثابت~$\ltrue$۔}{}
  \item د !!{propositional variable}s يو !!{denumerable}s سټ: $\Obj
    p_0$، $\Obj p_1$، $\Obj p_2$، \dots
  \item قضيوي نښلوونکي: \startycommalist
  \iftag{prvNot}{\ycomma $\lnot$ (نفي)}{}%
  \iftag{prvAnd}{\ycomma $\land$ (عطف)}{}%
  \iftag{prvOr}{\ycomma $\lor$ (فصل)}{}%
  \iftag{prvIf}{\ycomma $\lif$ (!!{conditional})}{}%
  \iftag{prvIff}{\ycomma $\liff$ (!!{biconditional})}{}.
  \item د تېر وخت عاملونه $\Ptemp$ او $\Htemp$۔
  \item د راتلونکي وخت عاملونه $\Ftemp$ او $\Gtemp$۔
\end{enumerate}
\end{defn}

وروسته به د مودالي عاملونو د نورو ډولونو د ورزياتولو امکان هم وڅېړو۔

\begin{defn}
د زماني ژبې \emph{!!^{formula}s} په استقرايي ډول داسې
  تعريفېږي:
\begin{enumerate}
\tagitem{prvFalse}{$\lfalse$ يو اټومي !!{formula} دے۔}{}

\tagitem{prvTrue}{$\ltrue$ يو اټومي !!{formula} دے۔}{}

\item هر قضيوي متغير~$\Obj p_i$ يو اټومي !!{formula} دے۔

\tagitem{prvNot}{که $!A$ يو !!a{formula} وي، نو $\lnot !A$ هم
  يو !!a{formula} دے۔}{}

\tagitem{prvAnd}{که $!A$ او $!B$ !!{formula}s وي، نو $(!A \land
  !B)$ يو !!a{formula} دے۔}{}

\tagitem{prvOr}{که $!A$ او $!B$ !!{formula}s وي، نو $(!A \lor !B)$
  يو !!a{formula} دے۔}{}

\tagitem{prvIf}{که $!A$ او $!B$ !!{formula}s وي، نو $(!A \lif !B)$
  يو !!a{formula} دے۔}{}

\tagitem{prvIff}{که $!A$ او $!B$ !!{formula}s وي، نو $(!A \liff !B)$
  يو !!a{formula} دے۔}{}

\item که $!A$ يو !!a{formula} وي، نو $\Ptemp  !A$، $\Htemp  !A$، $\Ftemp !A$، $\Gtemp  !A$ ټول
  !!{formula}s دي۔ په منجمده سرچينه کښې درېيم عامل په غلطه د عادي اېف
  په بڼه ليکل شوے؛ دلته د ژبې له پورته تعريف او د صدق له لاندې مادې سره
  سم د راتلونکي عامل رسمي نښه راوړل شوې ده۔

\tagitem{limitClause}{بل هېڅ شے !!a{formula} نۀ دے۔}{}
\end{enumerate}
\end{defn}

د نورو مفهومي منطقونو په څېر، د زماني منطقونو معنٰی پوهنه د اړيکيزو
مدلونو له مخې ورکول کېږي۔

\begin{defn}
  د زماني ژبې يو \emph{مدل} درې‌ګونې
  $\mModel{M} = \tuple{T, \prec, V}$ ده، چې په دې کښې
  \begin{enumerate}
  \item $T$ يو نا تش سټ دے، چې د وخت د ټکو په توګه تفسيرېږي۔
  \item $\prec$ پر~$T$ يوه دوه‌ځاييزه اړيکه ده۔
  \item $V$ داسې تابع ده چې هر !!{propositional
    variable}~$p$ ته د وخت د ټکو يو سټ $V(p)$ ټاکي۔
  \end{enumerate}
  که $t \prec t'$ وي، نو وايو چې $t$ له~$t'$ څخه \emph{مخکښې راځي}۔
  که $t \in V(p)$ وي، نو وايو چې $p$ په~$t$ کښې \emph{رښتيا دے}۔
\end{defn}

تر دې ځايه موږ د مخکښې‌والي پر اړيکه~$\prec$ هېڅ شرط نۀ دے ايښے۔
يعنې اوس د زماني مدلونو پر جوړښت هېڅ محدوديت نشته: وخت پکښې خطي،
څوڅانګيز، دوراني يا د بل هر ډول جوړښت لرلے شي۔

د عادي مودالي منطق په څېر، هر زماني مدل ټاکي چې کوم !!{formula}s
ئې په کومو ټکو کښې رښتيا دي۔ د اړيکيزې معنٰی پوهنې د بنسټيز مفهوم
دپاره هماغه نښه کاروو: «مدل~$\mModel{M}$ د ټکي~$t$ په حالت کښې
!!{formula}~$!A$ رښتيا کوي»۔ د رښتياوالي دا اړيکه په استقرايي ډول
تعريفېږي او د ټولو نا‌مودالي عاملونو دپاره له عادي مودالي حالت سره
يو شان ده۔

\begin{defn}\ollabel{defn:tmodels}
  په~$\mModel M$ کښې د~$t$ په حالت کښې \emph{د !!a{formula}~$!A$
  رښتياوالی}، چې په نښو کښې ئې $\mSat{M}{!A}[t]$ ليکو، په استقرايي
  ډول داسې تعريفېږي:
  \begin{enumerate}
  \tagitem{prvFalse}{\indcase{!A}{\lfalse}{$\mSat{M}{\lfalse}[t]$ هېڅکله نۀ وي}۔}{}
  \tagitem{prvTrue}{\indcase{!A}{\ltrue}{$\mSat{M}{\ltrue}[t]$ تل وي}۔}{}
  \item $\mSat{M}{p}[t]$ هله او يوازې هله چې $t \in V(p)$ وي۔
  \tagitem{prvNot}{\indcase{!A}{\lnot !B}{$\mSat{M}{\indfrm}[t]$ هله او يوازې هله چې
    $\mSat/{M}{!B}[t]$ وي}۔}{}
  \tagitem{prvAnd}{\indcase{!A}{(!B \land !C)}{$\mSat{M}{\indfrm}[t]$ هله او يوازې هله چې
    $\mSat{M}{!B}[t]$ او $\mSat{M}{!C}[t]$ دواړه وي}۔}{}
  \tagitem{prvOr}{\indcase{!A}{(!B \lor !C)}{$\mSat{M}{\indfrm}[t]$ هله او يوازې هله چې
    $\mSat{M}{!B}[t]$ يا $\mSat{M}{!C}[t]$ وي} (يا دواړه)۔}{}
  \tagitem{prvIf}{\indcase{!A}{(!B \lif !C)}{$\mSat{M}{\indfrm}[t]$ هله او يوازې هله چې
    $\mSat/{M}{!B}[t]$ يا $\mSat{M}{!C}[t]$ وي}۔}{}
  \tagitem{prvIff}{\indcase{!A}{(!B \liff !C)}{$\mSat{M}{\indfrm}[t]$ هله او يوازې هله چې
    يا $\mSat{M}{!B}[t]$ او $\mSat{M}{!C}[t]$ دواړه وي، يا
    نۀ $\mSat{M}{!B}[t]$ وي او نۀ $\mSat{M}{!C}[t]$}۔}{}
  \item\ollabel{defn:sub:mmodels-p}
    \indcase{!A}{\Ptemp  !B}{$\mSat{M}{\indfrm}[t]$ هله او يوازې هله چې
    د کوم $t' \in T$ دپاره چې $t' \prec t$ وي، $\mSat{M}{!B}[t']$ وي}
\item\ollabel{defn:sub:mmodels-h}
    \indcase{!A}{\Htemp  !B}{$\mSat{M}{\indfrm}[t]$ هله او يوازې هله چې
    د هر $t' \in T$ دپاره چې $t' \prec t$ وي، $\mSat{M}{!B}[t']$ وي}
   \item\ollabel{defn:sub:mmodels-f}
    \indcase{!A}{\Ftemp  !B}{$\mSat{M}{\indfrm}[t]$ هله او يوازې هله چې
    د کوم $t' \in T$ دپاره چې $t \prec t'$ وي، $\mSat{M}{!B}[t']$ وي}
\item\ollabel{defn:sub:mmodels-g}
    \indcase{!A}{\Gtemp  !B}{$\mSat{M}{\indfrm}[t]$ هله او يوازې هله چې
    د هر $t' \in T$ دپاره چې $t \prec t'$ وي، $\mSat{M}{!B}[t']$ وي}
  \end{enumerate}
\end{defn}

له دې معنٰی پوهنې څخه څرګندېږي چې $\Ptemp$ او~$\Htemp$ دوه‌ګوني
عاملونه دي؛ همداسې $\Ftemp$ او $\Gtemp$ هم دوه‌ګوني دي۔ نو
$\Htemp  !A$ د $\lnot \Ptemp \lnot !A$ په توګه تعريفولے شو،
او همدا خبره د $\Gtemp$ او~$\Ftemp$ په اړه هم کېږي۔

\end{document}
